%% XQdrant — ACM sigconf conversion (double-blind submission version)
%% Uses the official acmart class (sigconf, review, anonymous options).
%% Compile:
%%   pdflatex main && bibtex main && pdflatex main && pdflatex main
%%
%% BEFORE CAMERA-READY:
%%  - remove "review,anonymous" from \documentclass options
%%  - restore \author / \affiliation / \email
%%  - fill in \acmConference, \acmYear, \acmDOI, \acmISBN once ACM sends them
%%  - add \setcopyright{...} per ACM's rights email
%%  - re-point \texttt{...} artifact paths (see NOTE in Sec. 5.1 and Sec. 1)
%%    to an anonymized mirror (e.g. anonymous.4open.science) instead of a
%%    path that could identify the repository/author.

\documentclass[sigconf,review,anonymous]{acmart}
\usepackage{graphicx}
\graphicspath{{figures/}}

%% Suppress ACM's boilerplate rights/copyright block during review;
%% restore proper \setcopyright{...} + \acmConference{...} for camera-ready.
\setcopyright{none}
\acmConference[Anonymous Submission]{Anonymous Conference}{Anonymous Date}{Anonymous City, Country}
\acmYear{2026}
\acmDOI{}
\acmISBN{}
\acmPrice{}
\settopmatter{printacmref=false}
\renewcommand\footnotetextcopyrightpermission[1]{}

\begin{document}

\title{XQdrant: In-Database Attribution and Masked-Distance Subspace Search
  over HNSW}

%% Anonymous for double-blind review.
\author{Anonymous Author(s)}
\affiliation{%
  \institution{Anonymous Institution}
  \city{Anonymous City}
  \country{Anonymous Country}
}
\email{anonymous@domain.edu}

\renewcommand{\shortauthors}{Anonymous et al.}

\begin{abstract}
Vector databases return opaque scalar similarities over high-dimensional
embeddings, giving practitioners no coordinate-level explanation of a match and
no way to restrict HNSW traversal itself to a dynamic focus subspace. We present
\textbf{XQdrant}, an extension of the Qdrant engine addressing both problems, and
we report results for each honestly rather than only where they are favorable.

\emph{Attribution} is unambiguously positive: \texttt{with\_dims\_explained}
computes exact top-$m$ per-dimension score decomposition in-database, avoiding
client-side vector egress, and is recall-neutral against vanilla Qdrant while
running $\sim$2.2$\times$ faster than post-query extraction at encoder-scale
widths ($D\in\{768,1536\}$).

\emph{Masked-distance HNSW traversal} is a mixed result, and we treat that as
the paper's second contribution: a systems-level characterization of \emph{why}
subspace-restricted traversal is harder than it looks. Naive all-layer masking
collapses recall (SIFT1M Recall@10 as low as 0.006 at $D_{\mathrm{sub}}/D=0.1$)
because HNSW's graph edges encode full-space proximity. A hybrid design that
keeps full-distance routing on coarse layers and masks only near the base layer
(\texttt{mask\_from\_layer}) restores navigability (SIFT recall 0.79--0.98 for
cutoffs $\geq 1$) but does \emph{not} yield end-to-end latency wins, because
graph descent remains dominated by I/O and candidate-set bookkeeping rather than
distance arithmetic; a dedicated SIMD gather kernel wins microbenchmarks
(1.1--1.8$\times$ over repack-then-score) without changing this. We further show
that subspace/full-space $k$-NN \emph{coherence} predicts when a cheap
verification pass can recover full-space recall, and that an $\alpha$-weighted
blend of full and masked scores is strictly dominated by a hard layer cutoff.

We release a reproducible harness and an experiment index so each of these
positive and negative findings can be independently reproduced or cited.
\end{abstract}

\begin{CCSXML}
<ccs2012>
   <concept>
       <concept_id>10002951.10002952.10003190.10010509</concept_id>
       <concept_desc>Information systems~Nearest-neighbor search</concept_desc>
       <concept_significance>500</concept_significance>
   </concept>
   <concept>
       <concept_id>10002951.10002952.10003190</concept_id>
       <concept_desc>Information systems~Database query processing</concept_desc>
       <concept_significance>400</concept_significance>
   </concept>
   <concept>
       <concept_id>10010147.10010257.10010258.10010259.10010263</concept_id>
       <concept_desc>Computing methodologies~Feature selection</concept_desc>
       <concept_significance>300</concept_significance>
   </concept>
</ccs2012>
\end{CCSXML}

\ccsdesc[500]{Information systems~Nearest-neighbor search}
\ccsdesc[400]{Information systems~Database query processing}
\ccsdesc[300]{Computing methodologies~Feature selection}

\keywords{vector databases, HNSW, approximate nearest neighbor,
  explainability, feature attribution, subspace search, masked distance, SIMD}

\maketitle

%=============================================================================
\section{Introduction}
\label{sec:intro}

High-dimensional dense embeddings underpin retrieval-augmented generation
(RAG), recommendation, and semantic search~\cite{lewis2020rag,qdrant2024,wang2021milvus}.
Production engines traverse HNSW graphs~\cite{malkov2018hnsw}, collapsing every
comparison into a single scalar. This creates two gaps that motivate this paper.

\textbf{Explainability.} Practitioners cannot see which coordinates drove a hit
without fetching full vectors and recomputing per-dimension terms client-side —
an expensive, egress-bound pattern.

\textbf{Subspace routing.} Many workloads care about a dynamic focus set
$\mathcal{F}$ of coordinates (concept probes, ablations, modality masks), yet
stock HNSW always evaluates full $D$-dimensional distances during graph descent.
Rescoring on $\mathcal{F}$ after full-vector preselection improves semantic
specificity but, as we show, cannot reduce traversal cost.

\textbf{XQdrant} addresses the first gap successfully and investigates the
second rigorously. For attribution, \texttt{with\_dims\_explained} attaches
exact top-$m$ coordinate terms to search results in-database. For subspace
routing, \texttt{nearest.focus} with \texttt{masked=true} evaluates reduced
distances \emph{during} HNSW traversal — not just after it — via a hybrid layer
cutoff, gather-based SIMD kernels, an optional verify/overfetch pass, and a
full/masked score blend. Where the second contribution falls short of a
latency win, we do not paper over it: we identify the mechanism (navigability
vs.\ arithmetic cost) and report exactly which configurations are and are not
worth shipping.

\paragraph{Contributions.}
\begin{itemize}
  \item \textbf{In-database attribution (positive result).} Exact score
    decomposition with order-statistic top-$m$ selection; recall-neutral vs.\
    vanilla Qdrant and $\sim$2.2$\times$ faster than post-query extraction.
  \item \textbf{Masked-distance HNSW traversal (mechanism-level negative
    result).} A hybrid layer-cutoff design that restores navigability where
    naive masking collapses it, but does not translate into end-to-end
    speedup — and an explanation, grounded in coherence and kernel
    microbenchmarks, of why.
  \item \textbf{Coherence as a predictive diagnostic.} Subspace/full-space
    $k$-NN overlap predicts when a cheap verify pass can recover full-space
    recall and when it cannot, independent of the traversal strategy used.
  \item \textbf{Reproducible artifacts.} A curated experiment index and
    harness (anonymized link to be added after review; see
    Section~\ref{sec:eval}) mapping every claim above to a canonical script
    and dataset.
\end{itemize}

%=============================================================================
\section{Background}
\label{sec:background}

\subsection{HNSW and Contiguous SIMD}

HNSW maintains multi-layer proximity graphs~\cite{malkov2018hnsw}. Queries
greedily descend with candidate heap size \texttt{ef\_search}, then extract
top-$k$ at the base layer. Distance kernels assume contiguous \texttt{f32}
slices (AVX2/NEON), so a non-contiguous focus set must either (i) repack
selected coordinates into a scratch buffer before scoring, or (ii) gather them
directly via SIMD gather loads.

\subsection{Score Decomposition}

For Dot product, $s(\mathbf{q},\mathbf{v})=\sum_i q_i v_i$; Cosine applies the
same decomposition after $\ell_2$ normalization; Euclidean and Manhattan admit
analogous additive terms. Attribution is therefore exact accounting of the
engine's own metric, not a surrogate explanation model such as
SHAP~\cite{lundberg2017shap} or LIME~\cite{ribeiro2016lime}.

\subsection{Subspace Coherence}
\label{sec:coherence-def}

Let $N_k^{\mathrm{full}}(q)$ and $N_k^{\mathcal{F}}(q)$ denote exact $k$-NN sets
under the full metric and under focus set $\mathcal{F}$. We define
\emph{coherence} as $|N_k^{\mathrm{full}}\cap N_k^{\mathcal{F}}|/k$, and refer
to the diagnostic built around this quantity as \textbf{C1} in
Section~\ref{sec:eval-coherence}. Coherence is the quantity that determines,
independent of any particular traversal strategy, whether full-space neighbors
can even in principle be recovered from subspace-guided search — it recurs
throughout Section~\ref{sec:eval} as the explanatory variable behind several
otherwise-surprising results.

%=============================================================================
\section{In-Database Attribution}
\label{sec:attribution}

XQdrant injects attribution at the collection coordinator after shard merge,
with an early return when \texttt{with\_dims\_explained} is absent — zero
overhead on the unmodified hot path. Per explained hit, Algorithm~\ref{alg:attribution}
materializes $D$ terms and selects the top-$m$ by absolute contribution in
average $\mathcal{O}(D)$ time via \texttt{select\_nth\_unstable}.

\begin{figure}[htbp]
\centering
\begin{minipage}{0.92\linewidth}
\small
\textbf{Algorithm 1} (In-database top-$m$ attribution).\\
\textbf{Input:} query $\mathbf{q}$, vector $\mathbf{v}$, $D$, $m$, metric $\delta$.\\
1.\ Normalize operands if required by $\delta$.\\
2.\ For $i=0,\ldots,D-1$: append $(i,\textsc{Term}_i(\mathbf{q},\mathbf{v},\delta))$.\\
3.\ Return order-statistic top-$m$ by $|t|$, sorted descending.
\end{minipage}
\caption{In-database top-$m$ feature attribution.}
\label{alg:attribution}
\end{figure}

\texttt{nearest.focus} without masking still performs full HNSW preselection
and rescores on $\mathcal{F}$ afterward — useful for interpretability-constrained
ranking but, as Section~\ref{sec:eval-attr} shows, not a traversal-cost
reduction. That gap is exactly what motivates masked traversal
(Section~\ref{sec:masked}).

%=============================================================================
\section{Masked-Distance HNSW: Design Space}
\label{sec:masked}

We consider three traversal-level strategies for evaluating distance only over
a focus subspace $\mathcal{F}$ during HNSW descent, plus two supporting
mechanisms (a kernel choice and a recall-recovery pass) that apply across all
three. Table~\ref{tab:namemap} fixes the naming used for the rest of the paper.

\begin{table}[htbp]
  \centering
  \caption{Naming used throughout Sections~\ref{sec:masked}--\ref{sec:eval}.}
  \label{tab:namemap}
  \small
  \begin{tabular}{|l|p{5.4cm}|}
    \hline
    \textbf{Name} & \textbf{What it is} \\
    \hline
    M1 & Naive: mask every HNSW layer \\
    M2 & Hybrid: full distance for layers $\geq L$, masked below (\texttt{mask\_from\_layer}$=L$) \\
    M3 & M2 plus an $\alpha$-weighted blend of full/masked scores on masked layers \\
    \hline
    K1 & Kernel choice: repack-then-score vs.\ SIMD gather, orthogonal to M1--M3 \\
    V1 & Verify/overfetch: optional post-traversal full-distance rescore, orthogonal to M1--M3 \\
    C1 & Coherence diagnostic (Section~\ref{sec:coherence-def}), used to interpret M1--M3 and V1 \\
    \hline
  \end{tabular}
\end{table}

\subsection{API Surface}

Focus parameters ride inside the existing query body:
\begin{verbatim}
"focus": {
  "dims": [...], "masked": true,
  "mask_from_layer": 1,   // M2/M3 hybrid cutoff
  "verify": true,         // optional full rescore
  "alpha": 0.25           // M3 blend weight
}
\end{verbatim}
Constraints mirror existing dense local-search paths: no quantization/turbo, no
simultaneous \texttt{with\_dims\_explained} on masked paths, and fail-fast if
\texttt{mask\_from\_layer}, \texttt{verify}, or \texttt{alpha} are set without
\texttt{masked=true}.

\subsection{M1 --- Naive All-Layer Masking}

With \texttt{masked=true} and no cutoff, every HNSW comparison during descent
uses only coordinates in $\mathcal{F}$. This is the cheapest semantic change to
implement, and the one most likely to be tried first — which is precisely why
we report it as a cautionary baseline rather than a proposal: it destroys
navigability whenever the focus metric disagrees with the full-space geometry
the graph was built on (Section~\ref{sec:eval-m1}).

\subsection{M2 --- Hybrid Layer Cutoff}

Layer $\ell$ uses the masked scorer iff $\ell < L$ for cutoff
$L=\texttt{mask\_from\_layer}$; otherwise full distance is used. $L=0$ recovers
plain (unmasked) nearest search; a large $L$ recovers M1. The implementation
builds dual raw scorers inside \texttt{FilteredScorer} and switches via
\texttt{set\_layer} at layer boundaries — visited-set and heap logic are
unchanged, so this composes with the rest of the query pipeline for free.

\subsection{M3 --- Weighted Blend}

On masked layers, score $= \alpha \cdot s_{\mathrm{full}} + (1-\alpha)\cdot
s_{\mathrm{focus}}$; coarse (unmasked) layers are unaffected. Note that
blending computes \emph{both} distances, so it cannot reduce arithmetic cost on
its own — it exists purely as a routing-bias ablation to test whether it
recovers recall at a fixed cutoff more cheaply than raising $L$.

\subsection{K1 --- Kernel Choice: Repack vs.\ Gather}

Masked scoring can copy focus dims into a contiguous scratch buffer before
calling the existing SIMD kernels (\emph{repack}), or gather them directly
in-register (AVX2 \texttt{vgatherdps} for Dot/Cosine/Euclidean on x86, scalar
fallback elsewhere). We add the gather path, selectable via
\texttt{XQDRANT\_MASKED\_KERNEL}, and treat kernel choice as orthogonal to
M1--M3: any of the three can use either kernel.

\subsection{V1 --- Verify and Overfetch}

\texttt{verify=true} re-scores returned candidates with full distance after
traversal, reusing the existing quantization-rescore hook. At overfetch
$m{=}1$ (i.e., $\mathit{limit}=k$) this only reorders a size-$k$ set, so
Recall@$k$ is provably unchanged — it is a correctness/ordering feature, not a
recall-recovery one. Overfetch with $m>1$ ($\mathit{limit}=m\cdot k$, keep
top-$k$ after full rescore) turns masked traversal into a candidate
\emph{filter} and can recover recall, at a latency cost we quantify in
Section~\ref{sec:eval-verify}.

%=============================================================================
\section{Experimental Evaluation}
\label{sec:eval}

We structure this section as an argument rather than a list of runs: each
subsection's result motivates the next question. Table~\ref{tab:keep}
previews where this argument ends, so it can be read as a map before the
detailed results.

\subsection{Setup}

\paragraph{Attribution suite.} Synthetic corpora $N{=}10\mathrm{k}$, $Q{=}500$,
$D\in\{768,1536\}$, seed 42; baselines: vanilla Qdrant, post-query
extract-plus-client-top-$m$, and XQdrant. Run on an Apple Silicon laptop via
the suite orchestrator.

\paragraph{Masked-distance suite.} HTTP against a release XQdrant binary, with
\texttt{XQDRANT\_MASKED\_KERNEL=gather} unless a kernel comparison is the
subject of the experiment. Corpora: SIFT1M ($D{=}128$) and synthetic
$N{=}50\mathrm{k}$ at $D\in\{768,1536\}$, $Q{=}500$, $k{=}10$,
\texttt{ef\_search}$=128$. \emph{Absolute milliseconds are not compared across
the two suites, which run on different hosts; only within-suite comparisons
are load-bearing.} The full experiment index, scripts, and canonical CSVs for
every claim below will be released at an anonymized artifact link upon
publication (omitted here to preserve double-blind review); a
placeholder/anonymized mirror can be substituted before submission if the
venue requires an artifact URL at review time.

\subsection{Attribution and Focus Rescoring Are Not the Same Problem}
\label{sec:eval-attr}

Figure~\ref{fig:latency-recall} shows \texttt{with\_dims\_explained} preserves
the latency--recall frontier vs.\ vanilla Qdrant at $D{=}768$ (Recall@10 gap
$<0.01$). Figure~\ref{fig:attribution-depth} shows in-database attribution
holds $\sim$3.5\,ms p50 across $m\in\{1,\ldots,50\}$ while post-query
extraction sits near 7.8\,ms ($\approx 2.2\times$). This confirms attribution
as a clean systems win. But Figure~\ref{fig:subspace-rescore} shows that
\texttt{nearest.focus} \emph{rescoring} — the natural first attempt at
subspace search, using the same full-preselection machinery — yields speedup
$<1$ (0.84--0.87). Rescoring reduces the dimensionality of the arithmetic but
not the traversal it depends on. This negative result is what motivates
attacking the traversal itself, which is the subject of the rest of this
section.

\begin{figure}[htbp]
  \centering
  \includegraphics[width=0.92\linewidth]{plot1_latency_recall_d768.png}
  \caption{Attribution suite: latency--recall at $D{=}768$ with
    \texttt{with\_dims\_explained} ($m{=}10$).}
  \label{fig:latency-recall}
\end{figure}

\begin{figure}[htbp]
  \centering
  \includegraphics[width=0.92\linewidth]{plot3_attribution_depth_d768.png}
  \caption{Attribution depth: in-DB vs.\ post-query extraction ($D{=}768$).}
  \label{fig:attribution-depth}
\end{figure}

\begin{figure}[htbp]
  \centering
  \includegraphics[width=0.92\linewidth]{plot4_subspace_speedup_d768.png}
  \caption{Focus \emph{rescoring} speedup $T_{\mathrm{full}}/T_{\mathrm{sub}}$
    ($D{=}768$): values $<1$ show rescoring alone does not accelerate search.}
  \label{fig:subspace-rescore}
\end{figure}

\subsection{M1: Naive Masking Breaks Navigability}
\label{sec:eval-m1}

If rescoring cannot help, the natural next attempt is to mask distance
\emph{during} traversal (M1). Figure~\ref{fig:m1} shows this fails sharply on
SIFT1M: Recall@10 against full-space ground truth is approximately
0.006 / 0.079 / 0.257 / 0.464 / 0.989 at $D_{\mathrm{sub}}/D = 0.1 / 0.25 / 0.5
/ 0.75 / 1.0$. This is the navigability breakdown predicted in
Section~\ref{sec:background}: graph edges encode full-space proximity, and
masked greedy descent has no guarantee of reaching the region containing the
true neighbors once that proximity assumption is violated.

\begin{figure}[htbp]
  \centering
  \includegraphics[width=0.92\linewidth]{fig_option1_sift.png}
  \caption{M1 on SIFT1M: Recall@10 vs.\ full-space ground truth as
    $D_{\mathrm{sub}}/D$ shrinks.}
  \label{fig:m1}
\end{figure}

\subsection{M2: A Hybrid Cutoff Restores Navigability, Not Latency}
\label{sec:eval-m2}

M2 keeps full-distance routing on coarse layers, hypothesizing that most of
navigability is determined there. Figure~\ref{fig:m2} confirms this: on SIFT,
cutoffs $L\geq 1$ recover 0.79--0.98 recall (evaluated against subspace-restricted
ground truth), while $L{=}0$ tracks the low coherence floor. On high-$D$
synthetic data the same cutoff lifts otherwise-collapsed cells (e.g.,
$D{=}768$, ratio 0.25: $\sim$0.01 $\rightarrow$ $\sim$0.20) but plateaus near
0.15--0.22 — a ceiling we explain via coherence in
Section~\ref{sec:eval-coherence}.

Critically, \textbf{end-to-end speedup vs.\ full search remains below 1 in
every measured cell} (Figure~\ref{fig:m2-speed}), even using the gather kernel.
Section~\ref{sec:eval-kernel} isolates why: the arithmetic savings from masking
are real but small relative to the I/O and candidate-management cost that
dominates HNSW wall-clock time.

\begin{figure}[htbp]
  \centering
  \includegraphics[width=\linewidth]{fig_option2_sift_recall.png}
  \caption{M2 on SIFT1M: recall heatmap over \texttt{mask\_from\_layer}
    $\times$ subspace ratio, gather kernel.}
  \label{fig:m2}
\end{figure}

\begin{figure}[htbp]
  \centering
  \includegraphics[width=\linewidth]{fig_option2_d768_recall.png}
  \caption{M2 on synthetic $D{=}768$: same sweep, showing the recall
    plateau discussed in Section~\ref{sec:eval-coherence}.}
  \label{fig:m2-d768}
\end{figure}

\begin{figure}[htbp]
  \centering
  \includegraphics[width=0.92\linewidth]{fig_option2_sift_speedup.png}
  \caption{M2 SIFT end-to-end speedup vs.\ full search: no cell exceeds
    $1\times$, despite recovered recall.}
  \label{fig:m2-speed}
\end{figure}

\subsection{K1: Kernel Choice Explains the Ceiling, Not the Floor}
\label{sec:eval-kernel}

Figure~\ref{fig:kernel} isolates the SIMD kernel question directly: gather is
$\sim$1.24$\times$ / 1.14$\times$ / 1.42$\times$ / 1.77$\times$ faster than
repack-then-score at $D_{\mathrm{sub}}\in\{32,128,384,768\}$ in ns/op
microbenchmarks. Re-running the M2 sweep with gather instead of repack leaves
recall unchanged and still produces \textbf{zero} end-to-end cells above
$1\times$. This isolates the bottleneck to graph traversal overhead
(candidate heap management, visited-set bookkeeping, memory access to
non-contiguous graph pages) rather than the per-comparison arithmetic — the
kernel is not the limiting factor, so faster kernels alone cannot close the
gap identified in Section~\ref{sec:eval-m2}.

\begin{figure}[htbp]
  \centering
  \includegraphics[width=0.92\linewidth]{fig_xcut1_kernels.png}
  \caption{Masked Dot-product kernel throughput: repack vs.\ gather.}
  \label{fig:kernel}
\end{figure}

\subsection{C1: Coherence Explains the Recall Ceiling}
\label{sec:eval-coherence}

Why does M2's recall plateau at 0.15--0.22 on high-$D$ synthetic data
(Section~\ref{sec:eval-m2}) rather than approaching 1? Figure~\ref{fig:coherence}
measures subspace/full-space exact top-10 overlap directly: it grows with
ratio but stays low on synthetic $D{=}768$ ($\sim$0.007 $\rightarrow$ 0.30 from
ratio 0.1 $\rightarrow$ 0.75); SIFT is higher ($\sim$0.01 $\rightarrow$ 0.46).
This is a property of the \emph{data and focus set}, not of the traversal
algorithm: when the subspace and full-space neighbor sets barely overlap, no
amount of navigability recovery can produce full-space-correct results,
because the answers the two metrics agree on are simply different. This
reframes Section~\ref{sec:eval-m2}'s plateau from an implementation
limitation into an expected consequence of low coherence, and predicts when
the verification pass in Section~\ref{sec:eval-verify} can and cannot help.

\begin{figure}[htbp]
  \centering
  \includegraphics[width=\linewidth]{fig_xcut2_sift.png}
  \caption{Subspace$\leftrightarrow$full exact top-10 overlap vs.\ ratio, SIFT1M.}
  \label{fig:coherence}
\end{figure}

\begin{figure}[htbp]
  \centering
  \includegraphics[width=\linewidth]{fig_xcut2_d768.png}
  \caption{Same measurement on synthetic $D{=}768$: lower coherence throughout,
    consistent with the lower recall ceiling in Figure~\ref{fig:m2-d768}.}
  \label{fig:coherence-d768}
\end{figure}

\subsection{V1: Verify + Overfetch Trades Latency for Recall, Where Coherence Allows}
\label{sec:eval-verify}

Given that M2 cannot reach full-space recall on its own, can a post-traversal
verification pass close the gap? At overfetch $m{=}1$, masked and verified
Recall@10 are identical, as predicted in Section~\ref{sec:masked} — reordering
a fixed candidate set cannot change which candidates it contains. At $m{=}8$ on
SIFT with hybrid $L{=}1$ (high coherence regime), full-space recall rises
substantially: $0.26\rightarrow 0.66$ at ratio 0.5 and $0.45\rightarrow 0.89$ at
ratio 0.75 (Figure~\ref{fig:verify}), but at a cost of roughly $+45$\,ms p50
($\sim$1.5\,ms $\rightarrow$ $\sim$46\,ms) — overfetching by $8\times$ and
fully rescoring is not cheap. On high-$D$ synthetic data (low coherence),
recall stays $\leq \sim$0.19 even at $m{=}8$: exactly the outcome
Section~\ref{sec:eval-coherence} predicts, since no amount of overfetch
recovers neighbors that the subspace metric never surfaces as candidates in
the first place.

\begin{figure}[htbp]
  \centering
  \includegraphics[width=\linewidth]{fig_xcut3_sift_full.png}
  \caption{Full-space Recall@10, masked vs.\ verify, at overfetch $m\in\{1,8\}$,
    SIFT1M with hybrid $L{=}1$.}
  \label{fig:verify}
\end{figure}

\begin{figure}[htbp]
  \centering
  \includegraphics[width=\linewidth]{fig_xcut3_d768_full.png}
  \caption{Same measurement on synthetic $D{=}768$: overfetch cannot rescue
    recall under low coherence.}
  \label{fig:verify-d768}
\end{figure}

\subsection{M3: Blending Is Dominated by a Hard Cutoff}
\label{sec:eval-m3}

Finally, we test whether M3's continuous $\alpha$ can outperform M2's discrete
layer cutoff at matched cost. Figure~\ref{fig:m3} shows that on SIFT at ratio
0.25, recall is maximized at $\alpha{=}0$ for every cutoff tested
($\sim$0.84 for $L\geq 1$) and degrades monotonically toward the coherence
floor as $\alpha\rightarrow 1$. Since $\alpha{=}0$ is exactly M2, blending adds
API surface and computes both distance terms without ever improving on the
simpler design. We do not recommend shipping M3.

\begin{figure}[htbp]
  \centering
  \includegraphics[width=0.92\linewidth]{fig_option4_sift.png}
  \caption{M3 on SIFT: recall vs.\ $\alpha$ for cutoffs $\{0,1,2\}$ (ratio
    0.25, subspace ground truth).}
  \label{fig:m3}
\end{figure}

\subsection{Summary}

Table~\ref{tab:keep} consolidates the keep/skip verdict for every mechanism
evaluated above, following the argument chain: rescoring doesn't help
(\S\ref{sec:eval-attr}) $\to$ naive masking breaks navigability
(\S\ref{sec:eval-m1}) $\to$ a hybrid cutoff fixes navigability but not latency
(\S\ref{sec:eval-m2}) $\to$ the bottleneck is traversal overhead, not arithmetic
(\S\ref{sec:eval-kernel}) $\to$ the recall ceiling is a coherence property of the
data (\S\ref{sec:eval-coherence}) $\to$ verification can buy back recall only
where coherence allows, at real cost (\S\ref{sec:eval-verify}) $\to$ blending
adds nothing over the hard cutoff (\S\ref{sec:eval-m3}).

\begin{table}[htbp]
  \centering
  \caption{Keep / demote map for the masked-distance study.}
  \label{tab:keep}
  \small
  \begin{tabular}{|l|l|p{4.0cm}|}
    \hline
    \textbf{Mechanism} & \textbf{Verdict} & \textbf{Evidence} \\
    \hline
    Attribution (in-DB) & Keep & $\sim$2.2$\times$ vs.\ post-query; recall-neutral \\
    Focus rescoring only & Not accel. & Speedup $<1$ (\S\ref{sec:eval-attr}) \\
    M1 naive mask & Baseline only & Full-GT collapse on SIFT (\S\ref{sec:eval-m1}) \\
    M2 hybrid cutoff & \textbf{Keep} & Restores navigability; e2e speedup still $<1$ \\
    K1 gather kernel & Keep (micro) & 1.1--1.8$\times$ ns/op; does not close e2e gap \\
    C1 coherence & Keep & Predicts high-$D$ / verify limits \\
    V1, $m{=}1$ & API only & Recall@$k$ provably unchanged \\
    V1 + overfetch & Conditional & Helps on coherent data; $\sim$+45\,ms; fails on low coherence \\
    M3 blend & Skip & Best at $\alpha{=}0$, i.e.\ collapses to M2 \\
    \hline
  \end{tabular}
\end{table}

%=============================================================================
\section{Related Work}
\label{sec:related}

\paragraph{Vector databases.} Milvus~\cite{wang2021milvus},
Pinecone~\cite{pinecone2023}, and Qdrant~\cite{qdrant2024} expose ANN with
filters and quantization but not native per-coordinate decompositions or
layer-aware masked metrics.

\paragraph{ANN libraries.} FAISS~\cite{johnson2019faiss} and
ScaNN~\cite{guo2020scann} optimize large-scale search; explainability and
dynamic subspace routing remain post-processing concerns external to the
index.

\paragraph{Explainable AI.} SHAP and LIME~\cite{lundberg2017shap,ribeiro2016lime}
approximate model behavior via sampling or local surrogates. We instead exploit
additive metric structure for deterministic, zero-variance attributions aligned
with the engine's own ranking — closer to in-database OLAP
rollups~\cite{gray1998data} than to sampling-based XAI.

\paragraph{Dimension reduction and projected indexes.} Static projections or
multi-index layouts differ from \emph{query-time} focus sets evaluated against
an existing full-dimensional HNSW index. A dedicated auxiliary index built per
focus set would sidestep the navigability problem central to
Section~\ref{sec:eval-m1} entirely, at the cost of only covering a known,
registered palette of focus sets; we scoped this out of the current study and
discuss it as future work in Section~\ref{sec:discussion}.

%=============================================================================
\section{Discussion}
\label{sec:discussion}

\paragraph{What we recommend shipping.} In-database attribution and the M2
hybrid cutoff are the two positive systems results of this paper. The gather
kernel (K1) is worth keeping as an implementation detail regardless of
traversal strategy. Verify (V1) is useful as an optional correctness/recall-recovery
mode \emph{when combined with overfetch on coherent data}; it should not be
marketed as a general recall-recovery mechanism given Section~\ref{sec:eval-verify}'s
low-coherence failure mode. M3's $\alpha$-blend should not ship — it never
outperforms $\alpha{=}0$, which is simply M2.

\paragraph{Why masked traversal doesn't yield end-to-end speedup.} Even with a
faster masked-comparison kernel, HNSW still touches full vectors for storage
layout, graph structure, and — on hybrid coarse layers — full-precision SIMD.
Section~\ref{sec:eval-kernel}'s isolation experiment shows I/O and candidate
management dominate wall-clock time; the arithmetic XQdrant's masking saves is
simply not the bottleneck at the scales tested. This is a mechanism-level
finding, not an implementation shortcoming of our specific kernels, and we
believe it generalizes to any HNSW-family engine without a redesigned
traversal data structure (e.g., a masked-distance-aware graph built with a
dimension-aware edge structure from the start, rather than an existing
full-space graph reused at query time).

\paragraph{Auxiliary projected indexes (future work).} As noted in
Section~\ref{sec:related}, a secondary index built directly in a registered
focus subspace would have its own valid navigability guarantee and could
plausibly deliver the end-to-end speedup that query-time masking of an
existing full-space graph cannot. This is a natural next step but requires an
index-registration API and additional memory per registered focus set; we
leave a full implementation and evaluation to future work.

\paragraph{Limitations.} We did not build a dedicated high-$D$ archive for M1
(its collapse at high $D$ is evidenced indirectly via M2's low-$L$ cells and
the coherence measurements in Section~\ref{sec:eval-coherence}, not a direct
M1 sweep at $D\in\{768,1536\}$). We did not instrument visited-set divergence
between masked and full traversal directly, which would strengthen the
mechanistic story in Section~\ref{sec:eval-kernel}. The attribution suite and
masked-distance suite run on different hardware, so absolute latencies are not
compared across them, only within each suite. \textbf{All reported
recall/latency figures are currently single-run point estimates; repeated
trials with confidence intervals are needed before final submission and are
in progress.}

%=============================================================================
\section{Conclusion}
\label{sec:conclusion}

XQdrant shows that exact feature attribution can live inside a Rust vector
engine with negligible quality loss and a clear latency win over client-side
extraction. For masked-distance HNSW, we report a more nuanced result: a
hybrid layer cutoff is a genuine \emph{navigability} tool — the only one of
five traversal-level mechanisms we test that both keeps recall high and
requires no change to the underlying graph — but it is not yet a latency win,
and subspace/full-space coherence, not implementation effort, sets a hard
ceiling on how much recall a masked traversal of an unmodified full-space graph
can recover. Overfetch-and-verify and $\alpha$-blending clarify the remaining
design space: the former can recover recall when the data is coherent enough
to make masked candidates useful, at a real latency cost; the latter is
strictly dominated by a hard layer cutoff and should not ship. Together with a
public experiment index, these results give practitioners a concrete,
evidence-based answer to ``should I mask HNSW traversal by subspace,'' rather
than a single untested claim in either direction.

\begin{acks}
We thank the Qdrant open-source community for the engine baseline. (Full
acknowledgments to be restored in the camera-ready version, after review.)
\end{acks}

\bibliographystyle{ACM-Reference-Format}
\bibliography{main}

\end{document}
